\documentclass{amsart}
% For dealing with references we use the comment environment
\usepackage{verbatim}
\newenvironment{reference}{\comment}{\endcomment}
%\newenvironment{reference}{}{}
\newenvironment{slogan}{\comment}{\endcomment}
\newenvironment{history}{\comment}{\endcomment}

% For commutative diagrams we use Xy-pic
\usepackage[all]{xy}

% We use 2cell for 2-commutative diagrams.
\xyoption{2cell}
\UseAllTwocells

% We use multicol for the list of chapters between chapters
\usepackage{multicol}

% This is generally recommended for better output
\usepackage{lmodern}
\usepackage[T1]{fontenc}

% For cross-file-references

% Package for hypertext links:
\usepackage{hyperref}

% For any local file, say "hello.tex" you want to link to please

% Theorem environments.
%
\theoremstyle{plain}
\newtheorem{theorem}[subsection]{Theorem}
\newtheorem{proposition}[subsection]{Proposition}
\newtheorem{lemma}[subsection]{Lemma}

\theoremstyle{definition}
\newtheorem{definition}[subsection]{Definition}
\newtheorem{example}[subsection]{Example}
\newtheorem{exercise}[subsection]{Exercise}
\newtheorem{situation}[subsection]{Situation}

\theoremstyle{remark}
\newtheorem{remark}[subsection]{Remark}
\newtheorem{remarks}[subsection]{Remarks}

\numberwithin{equation}{subsection}

% Macros
%
\def\lim{\mathop{\mathrm{lim}}\nolimits}
\def\colim{\mathop{\mathrm{colim}}\nolimits}
\def\Spec{\mathop{\mathrm{Spec}}}
\def\Hom{\mathop{\mathrm{Hom}}\nolimits}
\def\Ext{\mathop{\mathrm{Ext}}\nolimits}
\def\SheafHom{\mathop{\mathcal{H}\!\mathit{om}}\nolimits}
\def\SheafExt{\mathop{\mathcal{E}\!\mathit{xt}}\nolimits}
\def\Sch{\mathit{Sch}}
\def\Mor{\mathop{\mathrm{Mor}}\nolimits}
\def\Ob{\mathop{\mathrm{Ob}}\nolimits}
\def\Sh{\mathop{\mathit{Sh}}\nolimits}
\def\NL{\mathop{N\!L}\nolimits}
\def\CH{\mathop{\mathrm{CH}}\nolimits}
\def\proetale{{pro\text{-}\acute{e}tale}}
\def\etale{{\acute{e}tale}}
\def\QCoh{\mathit{QCoh}}
\def\Ker{\mathop{\mathrm{Ker}}}
\def\Im{\mathop{\mathrm{Im}}}
\def\Coker{\mathop{\mathrm{Coker}}}
\def\Coim{\mathop{\mathrm{Coim}}}

% Boxtimes
%
\DeclareMathSymbol{\boxtimes}{\mathbin}{AMSa}{"02}

%
% Macros for moduli stacks/spaces
%
\def\QCohstack{\mathcal{QC}\!\mathit{oh}}
\def\Cohstack{\mathcal{C}\!\mathit{oh}}
\def\Spacesstack{\mathcal{S}\!\mathit{paces}}
\def\Quotfunctor{\mathrm{Quot}}
\def\Hilbfunctor{\mathrm{Hilb}}
\def\Curvesstack{\mathcal{C}\!\mathit{urves}}
\def\Polarizedstack{\mathcal{P}\!\mathit{olarized}}
\def\Complexesstack{\mathcal{C}\!\mathit{omplexes}}
% \Pic is the operator that assigns to X its picard group, usage \Pic(X)
% \Picardstack_{X/B} denotes the Picard stack of X over B
% \Picardfunctor_{X/B} denotes the Picard functor of X over B
\def\Pic{\mathop{\mathrm{Pic}}\nolimits}
\def\Picardstack{\mathcal{P}\!\mathit{ic}}
\def\Picardfunctor{\mathrm{Pic}}
\def\Deformationcategory{\mathcal{D}\!\mathit{ef}}

\setlength{\emergencystretch}{3em}
\begin{document}
\title[Stacks Project, Tag 03H0]{037 --- Stein factorization with geometrically connected fibres}
\maketitle
\noindent Copyright (C) 2005--2025 Johan de Jong. Stacks Project authors. Source: \href{https://stacks.math.columbia.edu/tag/03H0}{Tag 03H0}.

\noindent Editorial difficulty note (not author text). Difficulty H2: The relative spectrum of the direct image gives a finite factorization, but the proof must also establish geometrically connected fibres. The included normalization and proper-fibre arguments carry the nontrivial part of this claim..

\noindent Editorial provenance note (not author text). History: excerpt from revision \texttt{a0444\allowbreak 6e57e\allowbreak c1fbc\allowbreak 252a8\allowbreak 71afc\allowbreak ec775\allowbreak 2fb28\allowbreak 07b14}, more-morphisms.tex, lines 15204--15287. Reference presentation was adapted; original-author context and core support blocks were retained or added. The mathematical wording is unchanged. Licensed under GNU FDL 1.2 or later; no invariant sections or cover texts. See ../licenses/COPYING. Context blocks reproduce author definitions and supporting results; they are not additional counted problems.

\begin{definition}
\label{morphisms-definition-normalization-X-in-Y}
Let $f : Y \to X$ be a quasi-compact and quasi-separated morphism of schemes.
Let $\mathcal{O}'$ be the integral closure of $\mathcal{O}_X$ in
$f_*\mathcal{O}_Y$. The {\it normalization of $X$ in $Y$} is the
scheme\footnote{The scheme $X'$ need not be normal, for example if
$Y = X$ and $f = \text{id}_X$, then $X' = X$.}
$$
\nu : X' = \underline{\Spec}_X(\mathcal{O}') \to X
$$
over $X$. It comes equipped with a natural factorization
$$
Y \xrightarrow{f'} X' \xrightarrow{\nu} X
$$
of the initial morphism $f$.
\end{definition}

\begin{lemma}
\label{lemma-stein-universally-closed}
Let $S$ be a scheme. Let $f : X \to S$ be a universally closed and
quasi-separated morphism. There exists a factorization
$$
\xymatrix{
X \ar[rr]_{f'} \ar[rd]_f & & S' \ar[dl]^\pi \\
& S &
}
$$
with the following properties:
\begin{enumerate}
\item the morphism $f'$ is universally closed, quasi-compact, quasi-separated,
and surjective,
\item the morphism $\pi : S' \to S$ is integral,
\item we have $f'_*\mathcal{O}_X = \mathcal{O}_{S'}$,
\item we have $S' = \underline{\Spec}_S(f_*\mathcal{O}_X)$, and
\item $S'$ is the normalization of $S$ in $X$, see
Morphisms, Definition \ref{morphisms-definition-normalization-X-in-Y}.
\end{enumerate}
Formation of the factorization $f = \pi \circ f'$ commutes
with flat base change.
\end{lemma}

\begin{proof}
By Morphisms, Lemma \href{https://stacks.math.columbia.edu/tag/04XU}{lemma universally closed quasi compact (Tag 04XU)}
the morphism $f$ is quasi-compact. Hence the normalization $S'$ of $S$ in
$X$ is defined (Morphisms, Definition
\ref{morphisms-definition-normalization-X-in-Y})
and we have the factorization $X \to S' \to S$. By
Morphisms, Lemma \href{https://stacks.math.columbia.edu/tag/03GQ}{lemma normalization in universally closed (Tag 03GQ)}
we have (2), (4), and (5). The morphism $f'$ is universally closed by
Morphisms, Lemma \href{https://stacks.math.columbia.edu/tag/01W6}{lemma image proper scheme closed (Tag 01W6)}.
It is quasi-compact by
Schemes, Lemma \href{https://stacks.math.columbia.edu/tag/03GI}{lemma quasi compact permanence (Tag 03GI)}
and quasi-separated by
Schemes, Lemma \href{https://stacks.math.columbia.edu/tag/01KV}{lemma compose after separated (Tag 01KV)}.

\medskip\noindent
To show the remaining statements we may assume the base scheme $S$ is affine,
say $S = \Spec(R)$. Then $S' = \Spec(A)$ with
$A = \Gamma(X, \mathcal{O}_X)$ an integral $R$-algebra.
Thus it is clear that $f'_*\mathcal{O}_X$
is $\mathcal{O}_{S'}$ (because $f'_*\mathcal{O}_X$ is quasi-coherent,
by
Schemes, Lemma
\href{https://stacks.math.columbia.edu/tag/01LC}{lemma push forward quasi coherent (Tag 01LC)},
and hence equal to $\widetilde{A}$). This proves (3).

\medskip\noindent
Let us show that $f'$ is surjective. As $f'$ is universally closed (see above)
the image of $f'$ is a closed subset
$V(I) \subset S' = \Spec(A)$. Pick $h \in I$. Then
$h|_X = f^\sharp(h)$ is a global section of the structure sheaf of
$X$ which vanishes at every point. As $X$ is quasi-compact this means
that $h|_X$ is a nilpotent section, i.e., $h^n|X = 0$ for some $n > 0$.
But $A = \Gamma(X, \mathcal{O}_X)$, hence $h^n = 0$.
As every element of $I$ is nilpotent, we conclude that $V(I) = S'$ as desired.

\medskip\noindent
By Cohomology of Schemes, Lemma \href{https://stacks.math.columbia.edu/tag/02KH}{lemma flat base change cohomology (Tag 02KH)}
we see that formation of $f_*\mathcal{O}_X$ commutes with flat base change.
Formation of the relative spectrum commutes with any base change by
Constructions, Lemma \href{https://stacks.math.columbia.edu/tag/01LX}{lemma spec properties (Tag 01LX)}.
Thus formation of the factorization commutes with flat base change.
\end{proof}

\begin{theorem}[Stein factorization; Noetherian case]
\label{theorem-stein-factorization-Noetherian}
Let $S$ be a locally Noetherian scheme.
Let $f : X \to S$ be a proper morphism.
There exists a factorization
$$
\xymatrix{
X \ar[rr]_{f'} \ar[rd]_f & & S' \ar[dl]^\pi \\
& S &
}
$$
with the following properties:
\begin{enumerate}
\item the morphism $f'$ is proper with geometrically connected fibres,
\item the morphism $\pi : S' \to S$ is finite,
\item we have $f'_*\mathcal{O}_X = \mathcal{O}_{S'}$,
\item we have $S' = \underline{\Spec}_S(f_*\mathcal{O}_X)$, and
\item $S'$ is the normalization of $S$ in $X$, see
Morphisms, Definition \ref{morphisms-definition-normalization-X-in-Y}.
\end{enumerate}
\end{theorem}

\begin{proof}
Let $f = \pi \circ f'$ be the factorization of
Lemma \ref{lemma-stein-universally-closed}. Note that besides the
conclusions of Lemma \ref{lemma-stein-universally-closed} we
also have that $f'$ is separated
(Schemes, Lemma \href{https://stacks.math.columbia.edu/tag/01KV}{lemma compose after separated (Tag 01KV)})
and finite type
(Morphisms, Lemma \href{https://stacks.math.columbia.edu/tag/01T8}{lemma permanence finite type (Tag 01T8)}).
Hence $f'$ is proper. By
Cohomology of Schemes, Proposition
\href{https://stacks.math.columbia.edu/tag/02O5}{proposition proper pushforward coherent (Tag 02O5)}
we see that $f_*\mathcal{O}_X$ is a coherent $\mathcal{O}_S$-module.
Hence we see that $\pi$ is finite, i.e., (2) holds.

\medskip\noindent
This proves all but the most interesting assertion, namely that
all the fibres of $f'$ are geometrically connected.
It is clear from the discussion above that we may replace $S$ by $S'$,
and we may therefore assume that $S$ is Noetherian, affine,
$f : X \to S$ is proper, and $f_*\mathcal{O}_X = \mathcal{O}_S$.
Let $s \in S$ be a point of $S$. We have to show that $X_s$ is
geometrically connected. By Lemma
\href{https://stacks.math.columbia.edu/tag/03GZ}{lemma characterize geometrically connected fibres (Tag 03GZ)}
we see that it suffices to show $X_u$ is connected
for every \'etale neighbourhood $(U, u) \to (S, s)$.
We may assume $U$ is affine. Thus $U$ is Noetherian
(Morphisms, Lemma \href{https://stacks.math.columbia.edu/tag/01T6}{lemma finite type noetherian (Tag 01T6)}),
the base change $f_U : X_U \to U$ is proper
(Morphisms, Lemma \href{https://stacks.math.columbia.edu/tag/01W4}{lemma base change proper (Tag 01W4)}),
and that also $(f_U)_*\mathcal{O}_{X_U} = \mathcal{O}_U$
(Cohomology of Schemes, Lemma \href{https://stacks.math.columbia.edu/tag/02KH}{lemma flat base change cohomology (Tag 02KH)}).
Hence after replacing
$(f : X \to S, s)$ by the base change $(f_U : X_U \to U, u)$
it suffices to prove that the fibre $X_s$ is connected when
$f_*\mathcal{O}_X = \mathcal{O}_S$. We can deduce this
from Derived Categories of Schemes, Lemma
\href{https://stacks.math.columbia.edu/tag/0G7X}{lemma proper idempotent on fibre (Tag 0G7X)}
(by looking at idempotents in the structure sheaf of $X_s$)
but we will also give a direct argument below.

\medskip\noindent
Namely, we apply the theorem on formal functions,
more precisely Cohomology of Schemes, Lemma
\href{https://stacks.math.columbia.edu/tag/02OD}{lemma formal functions stalk (Tag 02OD)}.
It tells us that
$$
\mathcal{O}^\wedge_{S, s} = (f_*\mathcal{O}_X)_s^\wedge =
\lim_n H^0(X_n, \mathcal{O}_{X_n})
$$
where $X_n$ is the $n$th infinitesimal neighbourhood of $X_s$.
Since the underlying topological space of $X_n$ is equal to that
of $X_s$ we see that if $X_s = T_1 \amalg T_2$ is a disjoint union
of nonempty open and closed subschemes, then similarly
$X_n = T_{1, n} \amalg T_{2, n}$ for all $n$. And this in turn means
$H^0(X_n, \mathcal{O}_{X_n})$ contains a nontrivial idempotent $e_{1, n}$,
namely the function which is identically $1$ on $T_{1, n}$ and
identically $0$ on $T_{2, n}$. It is clear that $e_{1, n + 1}$
restricts to $e_{1, n}$ on $X_n$. Hence $e_1 = \lim e_{1, n}$
is a nontrivial idempotent of the limit. This contradicts the fact
that $\mathcal{O}^\wedge_{S, s}$ is a local ring. Thus the
assumption was wrong, i.e., $X_s$ is connected, and we win.
\end{proof}
\end{document}
